\section{Conclusion}
\label{sec:conclusion}

This project began with a straightforward question: could a high school student, working
independently, design and implement a functional thrust vector control system for a model rocket
from the ground up? The answer, after years of iterative hardware design, self-directed learning,
and rigorous analysis, is yes. The process of getting there is perhaps as significant as the
result itself.

From a technical standpoint, this paper has demonstrated that thrust vector control is a viable
stabilization mechanism at the hobby scale. The PID-Kalman control architecture, implemented on a
Teensy 4.1 microcontroller and driven by real-time MPU-6050 sensor data \cite{invensense_mpu6050, pjrc_teensy41}, produced stable closed-loop behavior across a range of simulated disturbance
scenarios. Simulation results consistently showed that the system could arrest angular deviations,
reject external disturbances, and operate within actuator limits without passive aerodynamic
stabilization. The Monte Carlo analysis further validated the robustness of the control
architecture under compounded uncertainty, demonstrating that failure modes were governed by
physical constraints rather than control instability.

It is important to be clear about the boundaries of these conclusions. The validation presented
here is entirely simulation-based. Airspace restrictions prevented powered flight testing, which
means that real-world phenomena---including structural vibration, aeroelastic coupling, launch
rail dynamics, and high-acceleration sensor behavior---remain unverified. The rigid-body assumption
underlying all simulations is a meaningful simplification, and the gap between simulated and actual
flight performance is something only a real launch can close. This paper should therefore be read
as a rigorous design and analysis study rather than a flight-validated system.

Several natural next steps follow from this work. A companion research study is currently underway
to directly quantify the accuracy of the simulation framework developed here, using a custom 2-DOF
static test stand to conduct live motor firings and compare measured angular error against
simulation predictions using the Integral of Absolute Error as the primary performance metric.
This work will identify where the simulation diverges from physical reality and trace that gap back
to specific unmodeled phenomena such as actuator rate saturation, gimbal friction, and vibration
noise. Beyond that, integrating a barometric pressure sensor would provide an independent altitude
measurement for more robust apogee detection. On the control side, adaptive gain scheduling---where
PID parameters shift as a function of thrust magnitude or estimated dynamic pressure---could
meaningfully improve performance during the thrust decay phase where control authority diminishes
most rapidly. Extending the platform to a multi-stage configuration would introduce additional
complexity in both the state machine and the control architecture, representing a logical
progression of the current design.

More broadly, the goal of this project was never exclusively technical. Thrust vector control
remains largely inaccessible in hobbyist and educational rocketry, not because the physics are
beyond reach, but because the path from concept to implementation is poorly documented at the
small scale. If this paper helps another student shortcut even a fraction of the trial and error
that went into this system, it will have done exactly what it was meant to do.
